\documentclass{article}
\usepackage[T1]{fontenc}
\usepackage{isabelle,isabellesym}
\usepackage{amsmath}
\usepackage{amssymb}

% this should be the last package used
\usepackage{pdfsetup}

% urls in roman style, theory text in math-similar italics
\urlstyle{rm}
\isabellestyle{it}

\begin{document}

\title{The Abel--Ruffini theorem and the unsolvability of the quintic}
\author{Lawrence C. Paulson}
\date{}

\maketitle

\begin{abstract}
  This entry proves the \emph{Abel--Ruffini theorem} --- a polynomial equation
  solvable by radicals has a solvable Galois group --- and applies it to a
  concrete quintic to conclude that the general quintic is not solvable by
  radicals.  We follow the corresponding \emph{mathlib} formalisation
  \cite{mathlib-abel-ruffini}.

  The forward direction --- \emph{solvable by radicals implies solvable Galois
  group} --- is assembled on the session
  \texttt{New\_Algebra}: cyclotomic and Kummer extensions have solvable (indeed
  abelian) Galois groups, solvability is preserved by the tower correspondence,
  and the normal closure of a radical tower is again a radical tower obtained by
  splicing the Galois conjugates of the original tower.

  For the concrete witness we study the rational quintic $\Phi = x^5 - 4x + 2$
  and show that the Galois group of its splitting field over $\mathbb{Q}$ is
  not solvable.  We prove that $\Phi$ is irreducible over $\mathbb{Q}$
  (Eisenstein's criterion at $2$) and, using the AFP
  entry \texttt{Sturm\_Sequences} \cite{eberl-sturm}, show that it has exactly three
  real roots and one conjugate pair of non-real roots. Complex conjugation
  therefore acts on the five roots as a transposition, while irreducibility
  makes the Galois action transitive; a transitive subgroup of $S_5$ containing
  a transposition is the whole of $S_5$, which is not solvable.  Combining the
  two halves yields the headline result that $\Phi$ is not solvable by radicals.
\end{abstract}

% sane default for proof documents
\parindent 0pt\parskip 0.5ex

% for proof reading
\pagestyle{plain}

\tableofcontents
\newpage

% generated text of all theories
\input{session}

% optional bibliography
\bibliographystyle{abbrv}
\bibliography{root}

\end{document}
